% GENERATED FILE. Edit canonical lessons, then run npm run generate:books.
% canonical-source-hash: fnv1a64:ef50d51d1e9dd9ba
% canonical-lessons: TE-C302-eddu, TE-C302-duda, TE-C302-gorre, TE-C302-todelu, TE-C302-mosali

\chapter{Ox, Calf, Sheep, Wolf, Crocodile}
\label{ch:te-ox-calf-sheep-wolf-crocodile}
\hlchaptermodality{302}

\begin{chapteropening}
\textbf{By the end of this chapter:} \emph{I can say an ox, a calf, a sheep, a wolf and a crocodile.}

Complete the last lesson of chapter 302: I can say an ox, a calf, a sheep, a wolf and a crocodile.
\end{chapteropening}

\section[\texorpdfstring{\te{ఎద్దు} (eddu)}{ఎద్దు (eddu)}]{\te{ఎద్దు} (eddu) — an ox, a bull}

\label{lesson:TE-C302-eddu}



\begin{quote}
Before the new one: say the Telugu for a taxi, then the Telugu for petrol.
\end{quote}

\subsection*{You'll want to know: \te{ఎద్దు}}
\textbf{\te{ఎద్దు}} — \emph{eddu} — \textquotedblleft{}an ox, a bull\textquotedblright{}.

\begin{cousinweb}
\textbf{\te{ఎద్దు}} (\emph{eddu}) is native Dravidian, one word with Tamil \emph{erutu} and Kannada \emph{eddu} (DEDR 815). DEDR lists Telugu \emph{eddu}, \textquotedblleft{}ox, bull\textquotedblright{}.
\end{cousinweb}

\begin{grammarlens}[title={What you've built}]
One more word for this chapter.
\end{grammarlens}

\subsection*{Your turn}
\begin{itemize}
\raggedright
  \item \emph{Say it:} \emph{eddu}
  \item \emph{Say it:} \emph{eddu}, once more
  \item \emph{Say it:} say \emph{eddu}
  \item \emph{Recall:} say the Telugu for a taxi, then the Telugu for petrol, then say \emph{eddu} again
\end{itemize}

\begin{tcolorbox}[breakable,colback=teal!4,colframe=teal!35!black,title={Before you move on}]
What does \te{ఎద్దు} mean? (\textquotedblleft{}An ox, a bull\textquotedblright{}.) Say it once more.
\end{tcolorbox}
\section[\texorpdfstring{\te{దూడ} (dūḍa)}{దూడ (dūḍa)}]{\te{దూడ} (dūḍa) — a calf}

\label{lesson:TE-C302-duda}



\begin{quote}
Before the new one: say the Telugu for petrol, then the Telugu for an ox.
\end{quote}

\subsection*{You'll want to know: \te{దూడ}}
\textbf{\te{దూడ}} — \emph{dūḍa} — \textquotedblleft{}a calf\textquotedblright{}.

\begin{cousinweb}
\textbf{\te{దూడ}} (\emph{dūḍa}) is Dravidian as far as DEDR records it: DEDR 3378 lists \emph{dūḍa}, \textquotedblleft{}a calf\textquotedblright{}.
\end{cousinweb}

\begin{grammarlens}[title={What you've built}]
One more word for this chapter.
\end{grammarlens}

\subsection*{Your turn}
\begin{itemize}
\raggedright
  \item \emph{Say it:} \emph{dūḍa}
  \item \emph{Say it:} \emph{dūḍa}, once more
  \item \emph{Say it:} say \emph{dūḍa}
  \item \emph{Recall:} say the Telugu for petrol, then the Telugu for an ox, then say \emph{dūḍa} again
\end{itemize}

\begin{tcolorbox}[breakable,colback=teal!4,colframe=teal!35!black,title={Before you move on}]
What does \te{దూడ} mean? (\textquotedblleft{}A calf\textquotedblright{}.) Say it once more.
\end{tcolorbox}
\section[\texorpdfstring{\te{గొర్రె} (gorre)}{గొర్రె (gorre)}]{\te{గొర్రె} (gorre) — a sheep}

\label{lesson:TE-C302-gorre}



\begin{quote}
Before the new one: say the Telugu for an ox, then the Telugu for a calf.
\end{quote}

\subsection*{You'll want to know: \te{గొర్రె}}
\textbf{\te{గొర్రె}} — \emph{gorre} — \textquotedblleft{}a sheep\textquotedblright{}.

\begin{cousinweb}
\textbf{\te{గొర్రె}} (\emph{gorre}) is native Dravidian, one word with Tamil \emph{koṟi} and Kannada \emph{koṟi} (DEDR 2165). DEDR lists Telugu \emph{gorre}, \textquotedblleft{}sheep\textquotedblright{}.
\end{cousinweb}

\begin{grammarlens}[title={What you've built}]
One more word for this chapter.
\end{grammarlens}

\subsection*{Your turn}
\begin{itemize}
\raggedright
  \item \emph{Say it:} \emph{gorre}
  \item \emph{Say it:} \emph{gorre}, once more
  \item \emph{Say it:} say \emph{gorre}
  \item \emph{Recall:} say the Telugu for an ox, then the Telugu for a calf, then say \emph{gorre} again
\end{itemize}

\begin{tcolorbox}[breakable,colback=teal!4,colframe=teal!35!black,title={Before you move on}]
What does \te{గొర్రె} mean? (\textquotedblleft{}A sheep\textquotedblright{}.) Say it once more.
\end{tcolorbox}
\section[\texorpdfstring{\te{తోడేలు} (tōḍēlu)}{తోడేలు (tōḍēlu)}]{\te{తోడేలు} (tōḍēlu) — a wolf}

\label{lesson:TE-C302-todelu}



\begin{quote}
Before the new one: say the Telugu for a calf, then the Telugu for a sheep.
\end{quote}

\subsection*{You'll want to know: \te{తోడేలు}}
\textbf{\te{తోడేలు}} — \emph{tōḍēlu} — \textquotedblleft{}a wolf\textquotedblright{}.

\begin{cousinweb}
\textbf{\te{తోడేలు}} (\emph{tōḍēlu}) is native Dravidian. DEDR 3548 lists a nasalised \emph{tōḍēlu}, \textquotedblleft{}wolf\textquotedblright{}, with Tamil \emph{tōṇṭāṉ} and Kannada \emph{tōḷa}, \textquotedblleft{}wolf\textquotedblright{}.
\end{cousinweb}

\begin{grammarlens}[title={What you've built}]
One more word for this chapter.
\end{grammarlens}

\subsection*{Your turn}
\begin{itemize}
\raggedright
  \item \emph{Say it:} \emph{tōḍēlu}
  \item \emph{Say it:} \emph{tōḍēlu}, once more
  \item \emph{Say it:} say \emph{tōḍēlu}
  \item \emph{Recall:} say the Telugu for a calf, then the Telugu for a sheep, then say \emph{tōḍēlu} again
\end{itemize}

\begin{tcolorbox}[breakable,colback=teal!4,colframe=teal!35!black,title={Before you move on}]
What does \te{తోడేలు} mean? (\textquotedblleft{}A wolf\textquotedblright{}.) Say it once more.
\end{tcolorbox}
\section[\texorpdfstring{\te{మొసలి} (mosali)}{మొసలి (mosali)}]{\te{మొసలి} (mosali) — a crocodile}

\label{lesson:TE-C302-mosali}



\begin{quote}
Before the new one: say the Telugu for a sheep, then the Telugu for a wolf.
\end{quote}

\subsection*{You'll want to know: \te{మొసలి}}
\textbf{\te{మొసలి}} — \emph{mosali} — \textquotedblleft{}a crocodile\textquotedblright{}.

\begin{cousinweb}
\textbf{\te{మొసలి}} (\emph{mosali}) is native Dravidian. DEDR 4952 lists \emph{mosali}, \textquotedblleft{}crocodile\textquotedblright{}, with Tamil \emph{mutalai} and Kannada \emph{mosaḷe}. DEDR compares Sanskrit \emph{mācala}, \textquotedblleft{}crocodile\textquotedblright{}, a lexicon word.
\end{cousinweb}

\begin{grammarlens}[title={What you've built}]
One more word for this chapter.
\end{grammarlens}

\subsection*{Your turn}
\begin{itemize}
\raggedright
  \item \emph{Say it:} \emph{mosali}
  \item \emph{Say it:} \emph{mosali}, once more
  \item \emph{Say it:} say \emph{mosali}
  \item \emph{Recall:} say the Telugu for a sheep, then the Telugu for a wolf, then say \emph{mosali} again
\end{itemize}

\begin{tcolorbox}[breakable,colback=teal!4,colframe=teal!35!black,title={Before you move on}]
What does \te{మొసలి} mean? (\textquotedblleft{}A crocodile\textquotedblright{}.) Say it once more.
\end{tcolorbox}
